\section*{الفصل \ref{ch:g11:ocean-ridges} --- الظهرات المحيطية والغلاف الصخري المحيطي}

\begin{solution}{exo:g11:ocean-ridges:1}
\omterm{def:g11:ocean-ridges:peridotite}{البيريدوتيت} صخر داكن خشن الحبيبات يتكون في معظمه من الأوليفين، مع شيء من
البيروكسين؛ ومنه يتكون جزء كبير من الوشاح. \omterm{def:g11:ocean-ridges:gabbro}{والغابرو} صخر ناري داكن خشن الحبيبات
من فلسبار البلاجيوكليز والبيروكسين، وهو أوفر صخور
\omterm{def:g11:ocean-ridges:oceanic-lithosphere}{القشرة المحيطية} العميقة. \omterm{def:g11:ocean-ridges:gabbro}{فالغابرو} ينتمي إلى القشرة، \omterm{def:g11:ocean-ridges:peridotite}{والبيريدوتيت} إلى الوشاح.
وقد تبلور \omterm{def:g11:ocean-ridges:gabbro}{الغابرو} ببطء، في أعماق تحت قاع البحر، حيث بردت الصهارة
ببطء: فأتيح لبلوراته الوقت لتنمو. أما لابة البازلت فبردت بسرعة على قاع
البحر، فبقيت بلوراتها صغيرة.
\end{solution}

\begin{solution}{exo:g11:ocean-ridges:2}
من قاع البحر نزولًا: الرواسب (غائبة عند المحور)، ولابات البازلت (وسائد
وصفائح منسابة)، والقواطع المتراصة، \omterm{def:g11:ocean-ridges:gabbro}{والغابرو}، ثم، تحت \omterm{def:g11:seismic-structure:moho}{الموهو}، \omterm{def:g11:ocean-ridges:peridotite}{بيريدوتيت}
الوشاح. وقد عبرت البئر $1507/7100 \approx 0.21$، أي نحو خُمس
القشرة: فلم يبلغ أي حفر قط \omterm{def:g11:ocean-ridges:gabbro}{الغابرو} العميق ولا \omterm{def:g11:seismic-structure:moho}{الموهو}.
\end{solution}

\begin{solution}{exo:g11:ocean-ridges:3}
$d(0) = \qty{2500}{m}$؛ $d(9) = 2500 + 350 \times 3 = \qty{3550}{m}$؛
$d(25) = 2500 + 350 \times 5 = \qty{4250}{m}$.
\end{solution}

\begin{solution}{exo:g11:ocean-ridges:4}
\omterm{def:g11:ocean-ridges:normal-fault}{الصدع العادي} صدع مائل تحركت على امتداده الكتلة الواقعة فوق مستوى
الصدع إلى الأسفل بالنسبة إلى الكتلة الواقعة تحته. وعند ظهرة بطيئة تشد الصفيحتان
\omterm{def:g11:lithosphere-asthenosphere:lithosphere}{الغلاف الصخري} البارد الهش في اتجاهين متعاكسين: فيتمطط وينكسر. وتنزلق
الكتل الواقعة فوق الصدوع إلى الأسفل، نحو المحور، والوادي المحوري هو
الكتلة التي هبطت بينها.
\end{solution}

\begin{solution}{exo:g11:ocean-ridges:5}
لا يغلي بسبب الضغط: فأكثر من كيلومترين من الماء
تضغط عليه، وتحت هذا الضغط يبقى الماء سائلًا في درجات أعلى بكثير من
\qty{100}{\celsius}. و«الدخان» الأسود سحابة من حبيبات دقيقة من كبريتيدات
الحديد ومعادن أخرى، تترسب حين يلتقي ماء الفتحة الساخن، الغني
بالفلزات الذائبة، بماء البحر البارد.
\end{solution}

\begin{solution}{exo:g11:ocean-ridges:6}
(1) على عمق نحو \qty{54}{km}، عند نحو \qty{1324}{\celsius}. (2) على عمق نحو
\qty{85}{km}، عند نحو $1400 + 0.45 \times 85 \approx
\qty{1438}{\celsius}$. (3) الوشاح الأسخن يبدأ بالانصهار أعمق، فينصهر
عمود أطول من الصخر الصاعد قبل أن يبلغ القشرة: فيصنع
بازلت أكثر، وقشرة أسمك.
\end{solution}

\begin{solution}{exo:g11:ocean-ridges:7}
(1) $1489 - 1345 = \qty{144}{\celsius}$ دون درجة انصهاره.
(2) يصعد $100 - 54 = \qty{46}{km}$؛ فتنغلق الفجوة بمقدار
$144/46 \approx \qty{3.1}{\celsius}$ لكل كيلومتر في المتوسط. (3) لا: ففي
طريقه إلى الأعلى يبرد الصخر قليلًا، بنحو $0.45 \times 46 \approx
\qty{21}{\celsius}$. إن درجة انصهاره هي التي تنخفض، وأسرع بكثير،
مع تناقص الضغط.
\end{solution}

\begin{solution}{exo:g11:ocean-ridges:8}
نصفا المعدلين \qty{7.3}{cm} و\qty{1.65}{cm} في السنة، أي
\qty{73}{km} و\qty{16.5}{km} في كل مليون سنة. ولصنع \qty{1000}{km} من
القاع على جانب واحد: $1000/73 \approx 13.7$ مليون سنة لظهرة شرق المحيط
الهادئ، و$1000/16.5 \approx 61$ مليون سنة لظهرة وسط الأطلسي.
\end{solution}

\begin{solution}{exo:g11:ocean-ridges:9}
(1) \qty{5}{cm} تساوي \qty{5e-5}{km}: $65000 \times 5 \times 10^{-5} =
\qty{3.25}{km^2}$ من القاع الجديد في السنة. (2) $3.25 \times 7.1 \approx
\qty{23}{km^3}$ من القشرة في السنة. (3) $3.6 \times 10^8 / 3.25 \approx 1.1
\times 10^8$ سنة، أي نحو 110 ملايين سنة: فبهذه الوتيرة يمكن للظهرات أن
تصنع قاع المحيطات كله في نحو مئة مليون سنة، وهي مدة قصيرة
مقارنة بعمر الأرض. ولهذا لا يوجد قاع محيطي قديم جدًا.
\end{solution}

\begin{solution}{exo:g11:ocean-ridges:10}
(1) يصنع مستوى الصدع زاوية \ang{60} مع الأفقي، والهبوط هو
الضلع الرأسي لمثلث قائم وتره الإزاحة على الصدع:
$1.5/\sin\ang{60} \approx \qty{1.73}{km}$. (2) يباعد كل صدع بين الجانبين
أفقيًا بمقدار $1.5/\tan\ang{60} \approx \qty{0.87}{km}$؛ والصدعان
معًا بنحو \qty{1.73}{km}. (3) لا. فمسافة \qty{25}{km} هي
البعد بين الصدعين اللذين يحدّان الكتلة الهابطة، لا مقدار
التمطط، الذي لا يتجاوز نحو \qty{1.7}{km}؛ والظهرة تواصل الانفتاح
بقواطع ولابات جديدة عند المحور، لا بالتصدع وحده.
\end{solution}

\begin{solution}{exo:g11:ocean-ridges:11}
$d(25) = 6400 - 3200\,e^{-25/62.8} \approx 6400 - 3200 \times 0.672 \approx
\qty{4251}{m}$؛ $d(100) = 6400 - 3200\,e^{-100/62.8} \approx 6400 - 3200
\times 0.203 \approx \qty{5749}{m}$. وبين 25 و100 مليون سنة
يغوص القاع بنحو \qty{1500}{m}. ويعطي قانون الجذر التربيعي
$2500 + 350 \times 5 = \qty{4250}{m}$ عند 25 مليون سنة: فالقانونان
متفقان حيث ينطبق كلاهما.
\end{solution}

\begin{solution}{exo:g11:ocean-ridges:12}
(1) $\alpha\,\Delta T\,L = 3.2 \times 10^{-5} \times 700 \times 100 =
\qty{2.24}{km}$. (2) $s = 3380 \times 2.24/(3380 - 1025) = 7571/2355 \approx
\qty{3.2}{km}$. (3) يسمح القانون الأسي بغوص كلي قدره $6400 - 2500 =
\qty{3900}{m}$: فالتبرد والانكماش وحدهما يفسران نحو
أربعة أخماسه، بقيم بسيطة لانخفاض درجة الحرارة وللسمك.
\end{solution}

\begin{solution}{exo:g11:ocean-ridges:13}
(1) $54 - 7 = \qty{47}{km}$ مقابل $85 - 7 = \qty{78}{km}$: فالعمود
الأسخن أطول بمقدار $78/47 \approx 1.66$ مرة. (2) $7.1 \times 1.66 \approx
\qty{11.8}{km}$. (3) السمك المقيس \qty{10.3}{km} أكبر من المعتاد، كما
كان متوقعًا، لكنه أقل من التقدير: فالكيلومترات الأعمق من العمود،
فوق بداية الانصهار مباشرة، تعطي مصهورًا أقل مما تعطيه الأقرب إلى السطح،
وافتراض القدر نفسه من المصهور لكل كيلومتر يبالغ في التقدير.
\end{solution}

\begin{solution}{exo:g11:ocean-ridges:14}
(1) يعطي القانون \qty{3178}{m} و\qty{4160}{m} و\qty{4782}{m}
و\qty{5267}{m}؛ والأوساط المقيسة أقل عمقًا بمقدار \qty{64}{m}
و\qty{142}{m} و\qty{317}{m} و\qty{619}{m}. (2) يكبر الفرق مع
العمر. (3) يعطي القانون عمق القاع نفسه، مصححًا من وزن
رواسبه؛ أما الخلايا المقيسة فتحمل رواسبها، الأسمك على
القاع الأقدم، وتشمل جبالًا بحرية وهضابًا بركانية، وهي
أقل عمقًا من القاع المحيط بها.
\end{solution}

\begin{solution}{exo:g11:ocean-ridges:15}
(1) $p = 1025 \times 9.8 \times 2500 \approx \qty{25.1}{MPa}$، مضافًا إليه
\qty{0.1}{MPa}: أي نحو \qty{25.2}{MPa}. (2) وهذا أعلى من الضغط
الحرج، \qty{22.064}{MPa}: فلا يمكن لماء الفتحة أن يغلي، مهما كانت حرارته.
(3) يبلغ الضغط \qty{22.064}{MPa} على العمق
$(22.064 - 0.1) \times 10^6/(1025 \times 9.8) \approx \qty{2190}{m}$: فكل
فتحة أعمق من نحو \qty{2200}{m} في مأمن من الغليان.
\end{solution}

\begin{solution}{pb:g11:ocean-ridges:1}
\textbf{1.} $14.6/2 = \qty{7.3}{cm}$ في السنة، أي \qty{73}{km} في كل مليون
سنة.

\textbf{2.} على عمق محور الظهرة، $d(0) = \qty{2500}{m}$.

\textbf{3.} من قاع البحر نزولًا: لابات البازلت (وسائد وصفائح منسابة)،
والقواطع المتراصة \omterm{def:g11:ocean-ridges:gabbro}{والغابرو}، أي نحو \qty{7.1}{km} من القشرة في المجموع؛ وتحت
\omterm{def:g11:seismic-structure:moho}{الموهو} \omterm{def:g11:ocean-ridges:peridotite}{بيريدوتيت} الوشاح، حتى قاعدة \omterm{def:g11:lithosphere-asthenosphere:lithosphere}{الغلاف الصخري}، التي لا تزال
قريبة جدًا من السطح عند المحور. وستستقر عليها الرواسب لاحقًا.

\textbf{4.} حين تتباعد الصفيحتان، يصعد \omterm{def:g11:lithosphere-asthenosphere:lithosphere}{الغلاف اللدن} تحت
الظهرة. ولا تنخفض درجة حرارته إلا بنحو \qty{0.45}{\celsius} لكل
كيلومتر، بينما تنخفض درجة انصهاره أسرع بكثير مع تناقص
الضغط. وعلى عمق نحو \qty{54}{km} تتساويان: فيبدأ الصخر بالانصهار،
ويستمر في الانصهار وهو يصعد، دون أن تُضاف إليه أي حرارة (\omterm{def:g11:ocean-ridges:partial-melting}{الانصهار
بتخفيف الضغط}).

\textbf{5.} في كل مليون سنة تصنع الظهرة $2 \times 73 = \qty{146}{km}$
من القاع على جانبيها؛ ولكل كيلومتر من الظهرة، يكون حجم القشرة
$146 \times 7.1 \times 1 \approx \qty{1040}{km^3}$.

\textbf{6.} $73 \times t$: \qty{73}{km} و\qty{292}{km} و\qty{1168}{km}
و\qty{4672}{km}.

\textbf{7.} $2000/73 \approx 27$ مليون سنة.

\textbf{8.} تغيّر الصفائح سرعتها واتجاهها على مدى عشرات الملايين من
السنين، كما يبيّن انعطاف سلسلة هاواي--الإمبراطور
(\cref{ch:g11:plate-motions})؛ فلا يلزم أن يبقى معدل اليوم ثابتًا طوال 64
مليون سنة.

\textbf{9.} بعيدًا عن المحور لا تعود القشرة فوق عدسة صهارة: فهي
تبرد، ولا تبقى تحتها حرارة كافية لتحريك الدوران وتسخين الماء.

\textbf{10.} غرب ظهرة شرق المحيط الهادئ تسافر القطعة على الصفيحة
الهادئية. ونحو الغرب تنتهي الصفيحة الهادئية على امتداد خنادق عميقة، حيث
تغوص تحت جاراتها: إنها نطاقات الاندساس (\cref{ch:g11:subduction-zones}).

\textbf{11.} $d(1) = \qty{2850}{m}$، $d(4) = 2500 + 700 = \qty{3200}{m}$،
$d(16) = 2500 + 1400 = \qty{3900}{m}$، $d(64) = 2500 + 2800 = \qty{5300}{m}$.

\textbf{12.} $3900 - 2500 = \qty{1400}{m}$ خلال ملايين السنين الستة عشر الأولى،
و$5300 - 3900 = \qty{1400}{m}$ بين 16 و64 مليون سنة: الغوص
نفسه، في زمن أطول بثلاث مرات. فالقاع يغوص أبطأ فأبطأ
كلما تقدم في العمر.

\textbf{13.} $350\sqrt{t} = 1500$ يعطي $\sqrt{t} \approx 4.29$، إذن
$t \approx 18.4$ مليون سنة.

\textbf{14.} $6400 - 3200\,e^{-64/62.8} \approx 6400 - 3200 \times 0.361
\approx \qty{5245}{m}$، أي أقل من قانون الجذر التربيعي بنحو \qty{55}{m}: فالقانونان
 متفقان في حدود نحو \qty{1}{\percent}.

\textbf{15.} الوسيط المقيس أقل عمقًا من القانون بنحو \qty{650}{m}.
فالقانون يصف القاع مصححًا من وزن رواسبه؛
أما الأعماق المقيسة فتشمل الرواسب، وهي على قاع كهذا حصيلة عشرات كثيرة من
ملايين السنين، وتشمل الجبال البحرية والهضاب البركانية، وكل ذلك يجعل
القاع المقيس أقل عمقًا.

\textbf{16.} $\alpha\,\Delta T = 3.2 \times 10^{-5} \times 700 = 0.0224$: فالجزء
الوشاحي أكثف بمقدار \qty{2.24}{\percent}،
$3380 \times 1.0224 \approx \qty{3456}{kg/m^3}$.

\textbf{17.} سمك الجزء الوشاحي $100 - 7.1 = \qty{92.9}{km}$. ومتوسط
الكثافة: $(7.1 \times 2900 + 92.9 \times 3456)/100 \approx (20590 + 321060)/100
\approx \qty{3416}{kg/m^3}$.

\textbf{18.} $3416 > 3380$: \omterm{def:g11:lithosphere-asthenosphere:lithosphere}{فالغلاف الصخري} القديم في مجمله أكثف من
\omterm{def:g11:lithosphere-asthenosphere:lithosphere}{الغلاف اللدن} الذي تحته بنحو \qty{36}{kg/m^3}، أي نحو \qty{1}{\percent}.
فلم تعد القشرة الخفيفة تكفي لتعويض جزئه الوشاحي البارد
الكثيف.

\textbf{19.} لأنه أكثف من \omterm{def:g11:lithosphere-asthenosphere:lithosphere}{الغلاف اللدن}، يمكنه أن يغوص عائدًا إلى
الوشاح عند خندق: إنه الاندساس، موضوع الفصل التالي.

\textbf{20.} عند 64 مليون سنة تقع قطعة القاع على عمق نحو
\qty{5300}{m}: أي \textbf{نحو 5.3~\unit{km}} تحت سطح البحر، أعمق بمقدار \qty{2.8}{km} من موضع
ولادتها.
\end{solution}
